% ============================================================
%  Chapter 6 – Performance Evaluation
% ============================================================
\chapter{Performance Evaluation}
\label{chap:perfeval}

\section{Introduction}

This chapter documents the testing methodology employed to evaluate WaveGuard V4 against the non-functional requirements specified in Chapter~\ref{chap:reqspec}.
Evaluation is structured around four principal dimensions: alert classification accuracy, system latency, API throughput, and alert threshold robustness.

\section{Testing Methodology}

\subsection{Simulated Sea-State Injection}

Because physical sea-state diversity cannot be reliably reproduced in a laboratory, WaveGuard V4 provides a dedicated simulation endpoint (\texttt{POST /api/test-buoy}) that injects synthetic telemetry directly into the backend without requiring physical hardware.
This allows deterministic, reproducible testing of the classification logic across the full range of alert states.

The test procedure was as follows:
\begin{enumerate}
  \item Define a test matrix of 150 scenario pairs \((m_i, h_i)\) spanning the three alert regions.
  \item Inject each pair via \texttt{curl} into \texttt{/api/test-buoy}.
  \item Poll \texttt{/api/status} and record the returned \texttt{alert} field.
  \item Compare against the ground truth derived from Equation~\ref{eq:fusion}.
\end{enumerate}

\noindent An example injection command is:
\begin{verbatim}
curl -X POST "http://localhost:8005/api/test-buoy?motion=0.12&
              device_status=HIGH%20WAVE"
\end{verbatim}

\subsection{Load Testing}

API throughput was measured using \texttt{wrk} — a modern HTTP benchmarking tool.
The \texttt{/api/status} endpoint was subjected to a 30-second load test at increasing concurrency levels (1, 10, 25, 50, 100 threads).

\subsection{Latency Measurement}

End-to-end latency was measured from the moment the ESP32 transmitted a buoy reading to the moment the updated alert badge appeared on the frontend browser.
A stopwatch timer was implemented in JavaScript (\texttt{performance.now()}) and a microcontroller timestamp was embedded in the HTTP request payload.

\section{Alert Classification Accuracy}

\subsection{Test Matrix Design}

The 150 test scenarios were distributed as follows:
\begin{itemize}
  \item 50 scenarios in the NORMAL region (\(\bar{a} \in [0, 0.04]\) g, \(H_s \in [0, 1.4]\) m)
  \item 50 scenarios in the WATCH region (\(\bar{a} \in [0.05, 0.09]\) g, \(H_s \in [1.5, 2.4]\) m)
  \item 50 scenarios in the WARNING region (\(\bar{a} \in [0.10, 0.50]\) g, \(H_s \in [2.5, 6.0]\) m)
\end{itemize}

\subsection{Confusion Matrix}

Table~\ref{tab:confusion} presents the classification confusion matrix from 150 test scenarios.

\begin{table}[H]
\centering
\caption{Alert Classification Confusion Matrix (\(n = 150\))}
\label{tab:confusion}
\begin{tabular}{lcccc}
\toprule
& \multicolumn{3}{c}{\textbf{Predicted}} & \\
\cmidrule(r){2-4}
\textbf{Actual} & NORMAL & WATCH & WARNING & \textbf{Total} \\
\midrule
NORMAL  & 50 & 0 & 0 & 50 \\
WATCH   & 0  & 50 & 0 & 50 \\
WARNING & 0  & 0  & 50 & 50 \\
\midrule
\textbf{Total} & 50 & 50 & 50 & \textbf{150} \\
\bottomrule
\end{tabular}
\end{table}

\noindent All 150 injected scenarios were classified correctly, yielding:
\begin{align}
  \text{Overall Accuracy} &= \frac{150}{150} = 100\% \label{eq:acc} \\
  \text{Precision (per class)} &= 1.00 \\
  \text{Recall (per class)}    &= 1.00 \\
  F_1 \text{ score}            &= 1.00
\end{align}

These results are expected for a rule-based classifier with deterministic thresholds; they confirm correct implementation of the fusion matrix with no boundary condition errors.

\subsection{Boundary Condition Testing}

Additional boundary tests were performed at the exact classification thresholds to verify correct inequality handling:

\begin{table}[H]
\centering
\caption{Boundary Condition Test Results}
\label{tab:boundary}
\begin{tabular}{cclc}
\toprule
\textbf{Input \(\bar{a}\) (g)} & \textbf{Input \(H_s\) (m)} & \textbf{Expected} & \textbf{Result} \\
\midrule
0.049 & 1.49 & NORMAL  & NORMAL  \\
0.050 & 1.50 & WATCH   & WATCH   \\
0.099 & 2.49 & WATCH   & WATCH   \\
0.100 & 2.50 & WARNING & WARNING \\
0.100 & 2.49 & WATCH   & WATCH   \\
0.099 & 2.50 & WATCH   & WATCH   \\
\bottomrule
\end{tabular}
\end{table}

All six boundary conditions produced the expected output, confirming that the threshold comparisons use non-strict inequalities as specified in Equation~\ref{eq:fusion}.

\section{System Latency}

\subsection{End-to-End Alert Latency}

Table~\ref{tab:latency} reports the measured end-to-end latency components.

\begin{table}[H]
\centering
\caption{End-to-End Alert Latency Components}
\label{tab:latency}
\begin{tabular}{lccc}
\toprule
\textbf{Component} & \textbf{Min (ms)} & \textbf{Mean (ms)} & \textbf{Max (ms)} \\
\midrule
ESP32 HTTP POST transmission & 12 & 18 & 45 \\
FastAPI endpoint processing  & 2  & 4  & 12 \\
SQLite write commit          & 1  & 2  & 8  \\
Open-Meteo API fetch (uncached) & 180 & 290 & 680 \\
Open-Meteo cache hit          & 0  & 0  & 1  \\
Frontend poll interval        & 0  & 2500 & 5000 \\
React re-render               & 3  & 5  & 15 \\
\midrule
\textbf{Total (cache miss)}   & 198 & 2819 & 6760 \\
\textbf{Total (cache hit)}    & 18  & 2529 & 5081 \\
\bottomrule
\end{tabular}
\end{table}

The dominant latency contributor is the 5-second frontend polling interval, not the backend processing time.
The mean end-to-end latency of approximately 2.5 seconds satisfies NFR-P2 (maximum 10 seconds).

\subsection{API Response Time Distribution}

Under a simulated load of 50 concurrent clients polling \texttt{/api/status} for 30 seconds:
\begin{itemize}
  \item P50 response time: 6 ms
  \item P95 response time: 18 ms
  \item P99 response time: 34 ms
  \item Maximum response time: 87 ms
  \item Throughput: 2,840 requests/second
\end{itemize}

These results comfortably satisfy NFR-P1 (500 ms target) and NFR-P3 (100 concurrent connections).

\section{Security Testing}

\subsection{Rate Limiter Verification}

Six successive login attempts were submitted from the same IP within one second.
Attempts 1–5 returned HTTP 200 (success) or HTTP 401 (wrong credentials).
Attempt 6 returned HTTP 429 (Too Many Requests), confirming correct rate limiter activation.

\subsection{Token Expiry}

A valid session token was extracted and used to access a protected admin endpoint immediately after login (successful, HTTP 200).
The same token was replayed after 28,800 seconds (8 hours) and received HTTP 401, confirming token TTL enforcement.

\section{Reliability Testing}

The system was left running continuously for 96 hours in hardware-present mode (ESP32 transmitting every second).
Across this period:
\begin{itemize}
  \item Total buoy readings recorded: 345,600
  \item Readings lost (network timeout): 42 (0.012\%)
  \item Backend crashes: 0
  \item Database corruption events: 0
\end{itemize}

The 99.988\% data capture rate satisfies NFR-R1 (72-hour continuous operation) and NFR-R2 (data integrity).

\section{Summary}

Performance evaluation demonstrates that WaveGuard V4 meets all quantitative non-functional requirements.
The rule-based alert classifier achieves 100\% accuracy on deterministic test inputs, end-to-end latency remains well below the 10-second target, API throughput exceeds requirements by more than 28×, and the system demonstrated high reliability over a 96-hour soak test.
